\documentclass[a4paper,11pt]{scrartcl}

\usepackage[T1]{fontenc}
\usepackage[utf8]{inputenc}
\usepackage[ngerman]{babel}
\usepackage{lmodern}
\usepackage{amsmath}
\usepackage{amssymb}
\usepackage[a4paper,margin=2.5cm]{geometry}
\usepackage{enumitem}
\usepackage{booktabs}
\usepackage[hidelinks]{hyperref}
\usepackage{catppuccin-dark}

\newcounter{aufgabe}
\newcommand{\Aufgabe}[2]{%
  \refstepcounter{aufgabe}%
  \subsection*{Aufgabe~\theaufgabe: #1 \hfill (#2~Punkte)}%
  \label{auf:\theaufgabe}}

\title{Übungsblatt 6: Numerical Integration (Numerische Integration)}
\subtitle{Begleitmaterial zur Vorlesung \glqq Numerical Methods\grqq{} (Prof.\ Dr.-Ing.\ Mark Schutera)}
\author{Quadraturregeln, zusammengesetzte Regeln und Monte-Carlo-Integration\\(Skript S.~57--70)}
\date{\today}

\begin{document}

\maketitle

\noindent\textbf{Gesamt: 57 Punkte.}\quad
Hilfsmittel: Taschenrechner bzw.\ die angegebenen Wertetabellen. Runden Sie, sofern nicht anders angegeben, auf sechs Nachkommastellen. Die Aufgaben steigen im Schwierigkeitsgrad an: Aufgaben 1--3 prüfen das Begriffsverständnis, Aufgaben 4--6 sind Rechenaufgaben analog zum Skript (mit eigenem Beispiel), Aufgaben 7--8 sind Transferaufgaben. Die Gleichungsnummern (42)--(51) und die Tabelle~3 beziehen sich auf das Skript.

\bigskip
\hrule

\section*{Aufgaben}

\Aufgabe{Glättung durch Integration (Smoothing) und der Mittelungseffekt}{5}

Das Skript motiviert die numerische Integration (Gl.~(42), $I = \int_a^b f(x)\,dx$) über das Problem hochfrequenter Loss-Landschaften: Eine Funktion mit Tausenden winziger, scharfer lokaler Minima --- wie eine Hochfrequenz-Version von Shekel's Foxholes --- lässt Gradient Descent sofort im ersten mikroskopischen \glqq Schlagloch\grqq{} (pothole) stecken bleiben. Als Beispiel dient $f(x) = \sin(x) + \sin(2\pi x)$ auf $[-3, 1]$.

\begin{enumerate}[label=(\alph*)]
  \item Erklären Sie, warum die Ersetzung der Punktauswertung $f(x)$ durch das (durch die Intervallbreite $h$ geteilte) Integral $\frac{1}{h}\int_{x-h/2}^{x+h/2} f(t)\,dt$ die Loss-Landschaft glättet. Nennen Sie die beiden \glqq guten Gründe\grqq{} aus dem Skript, numerische Integration in der Optimierung einzusetzen. \hfill (2~P)
  \item \emph{What-if-Frage des Skripts (S.~59):} Was passiert, wenn man über ein Intervall integriert, das zu destruktiver Interferenz (destructive interference) der höherfrequenten Sinuskomponente führt? \glqq Think wavelengths\grqq: Welche Intervallbreite $h$ würde den Anteil $\sin(2\pi x)$ sogar \emph{exakt} eliminieren, und warum wirkt die im Skript gewählte Breite $h = 0.48$ bereits fast genauso gut? \hfill (2~P)
  \item In welchem Sinne ist die Mittelpunktsregel (midpoint rule) ein Spezialfall der Monte-Carlo-Schätzung? \hfill (1~P)
\end{enumerate}

\schreibraum[8]
\Aufgabe{Midpoint, Trapez und Simpson konzeptionell --- und wann man aufhört}{6}

\begin{enumerate}[label=(\alph*)]
  \item Ordnen Sie den drei Quadraturregeln --- Mittelpunktsregel (midpoint rule), Trapezregel (trapezoid rule), Simpson-Regel (Simpson's rule) --- jeweils das geometrische Objekt bzw.\ den Polynomgrad zu, mit dem die Funktion auf einem Teilintervall approximiert wird. Warum liefert die Mittelpunktsregel intuitiv bessere Ergebnisse als die Links- oder Rechts-Endpunkt-Regeln (left/right endpoint rules)? \hfill (2~P)
  \item Ordnen Sie den vier Methoden Left/Right Endpoint, Midpoint, Trapezoid und Simpson's 1/3 die folgenden (gemischten) Fehlerordnungen aus Tabelle~3 des Skripts zu --- jeweils Einzelintervallfehler (single interval error) und akkumulierter Fehler (accumulated error):
  \[ O(h),\quad O(h^2),\quad O(h^2),\quad O(h^2),\quad O(h^3),\quad O(h^3),\quad O(h^4),\quad O(h^5). \]
  \hfill (2~P)
  \item \emph{Practical advice des Skripts: \glqq Stop at Simpson.\grqq} Warum sollte man für höhere Genauigkeit zusammengesetzte Regeln (composite rules) verwenden, statt den Polynomgrad der Interpolation immer weiter zu erhöhen? Erläutern Sie dazu das Runge-Phänomen (Runge's phenomenon) und das Problem hoher Grade $n \geq 8$. \hfill (2~P)
\end{enumerate}

\schreibraum[9]
\Aufgabe{Fluch der Dimensionalität (Curse of Dimensionality)}{5}

Für $d$-dimensionale Integrale benötigt deterministische Quadratur mit $N$ Stützstellen pro Dimension insgesamt $N^d$ Punkte.

\begin{enumerate}[label=(\alph*)]
  \item Begründen Sie kurz, warum $N^d$ Punkte nötig sind. Berechnen Sie die Gesamtzahl der Gitterpunkte für $N = 100$ Punkte pro Dimension bei $d = 2$ und bei $d = 10$. Wie lange bräuchte ein Rechner, der $10^9$ Funktionsauswertungen pro Sekunde schafft, für den Fall $d = 10$ (Größenordnung in Jahren)? \hfill (2~P)
  \item Das Skript betrachtet ein neuronales Netz mit $d = 10^6$ Parametern. Wie viele Gitterpunkte ergäbe schon das gröbstmögliche Gitter $N = 2$? Schätzen Sie die Anzahl der Dezimalstellen dieser Zahl ab (Hinweis: $\log_{10} 2 \approx 0.30103$). \hfill (2~P)
  \item Welche Eigenschaft der Monte-Carlo-Integration macht sie in hohen Dimensionen attraktiv, obwohl Quadraturregeln in den Konvergenz\emph{raten} überlegen sind? \hfill (1~P)
\end{enumerate}

\subsection*{Gemeinsames Szenario für die Aufgaben 4--7}

Zu berechnen ist das Integral
\[ I = \int_0^2 e^x\,dx. \]
Der exakte Wert dient jeweils zur Fehlerquantifizierung. \textbf{Wertetabelle} (auf sechs Nachkommastellen gerundet):
\begin{center}
\begin{tabular}{lccccc}
\toprule
$x$ & $0$ & $0.5$ & $1$ & $1.5$ & $2$ \\
$e^x$ & $1.000000$ & $1.648721$ & $2.718282$ & $4.481689$ & $7.389056$ \\
\midrule
$x$ & $0.25$ & $0.75$ & $1.75$ & & \\
$e^x$ & $1.284025$ & $2.117000$ & $5.754603$ & & \\
\bottomrule
\end{tabular}
\end{center}

\schreibraum[8]
\Aufgabe{Midpoint, Trapez und Simpson von Hand ($n = 1$ Intervall)}{8}

\begin{enumerate}[label=(\alph*)]
  \item Berechnen Sie den exakten Wert $\tilde{I} = \int_0^2 e^x\,dx$ analytisch. \hfill (1~P)
  \item Approximieren Sie $I$ mit der Mittelpunktsregel mit $n = 1$ Intervall, $\hat{I} = (b-a)\cdot f(m)$, und geben Sie den absoluten Fehler $|\tilde{I} - \hat{I}|$ an. \hfill (2~P)
  \item Approximieren Sie $I$ mit der Trapezregel (Gl.~(44)), $\int_a^b f(x)\,dx \approx \frac{h}{2}[f(a) + f(b)]$ mit $h = b - a$, inklusive Fehler. \hfill (2~P)
  \item Approximieren Sie $I$ mit der Simpson-Regel auf dem Gesamtintervall, $\int_a^b f(x)\,dx \approx \frac{h}{6}[f(a) + 4f(m) + f(b)]$ mit $h = b - a$ und Mittelpunkt $m$, inklusive Fehler. \hfill (2~P)
  \item Ordnen Sie die drei Fehler der Größe nach. Die Funktion $e^x$ ist konvex --- erklären Sie damit, warum die Mittelpunktsregel das Integral \emph{unterschätzt} und die Trapezregel es \emph{überschätzt}. \hfill (1~P)
\end{enumerate}

\schreibraum[11]
\Aufgabe{Zusammengesetzte Regeln (Composite Rules)}{8}

\begin{enumerate}[label=(\alph*)]
  \item Stellen Sie die zusammengesetzte Trapezregel und die zusammengesetzte Simpson-Regel für äquidistante Gitter $x_i = a + ih$, $i = 0, \ldots, n$ ($n$ gerade bei Simpson) allgemein auf. Geben Sie jeweils das Gewichtsmuster an und erklären Sie bei der Trapezregel, warum innere Punkte das Gewicht 2 tragen. \emph{Hinweis:} Bei der Composite-Simpson-Regel beträgt der Vorfaktor pro Teilintervall-\emph{Paar} der Breite $2h$ gerade $\frac{2h}{6} = \frac{h}{3}$ [Korrektur ggü.\ Original, S.~64: Die Zwischenzeilen der Herleitung im Skript drucken fälschlich $\frac{h}{6}$ pro Paar; die Endformel mit $\frac{h}{3}$ ist korrekt]. \hfill (2~P)
  \item Berechnen Sie $\int_0^2 e^x\,dx$ mit der zusammengesetzten Mittelpunktsregel mit $n = 2$ Intervallen ($h = 1$, Mittelpunkte $0.5$ und $1.5$), inklusive Fehler. \hfill (2~P)
  \item Berechnen Sie das Integral mit der zusammengesetzten Trapezregel mit $n = 2$ Intervallen ($h = 1$), inklusive Fehler. \hfill (2~P)
  \item Berechnen Sie das Integral mit der zusammengesetzten Simpson-Regel mit $n = 4$ Intervallen ($h = 0.5$), inklusive Fehler. \hfill (2~P)
\end{enumerate}

\schreibraum[11]
\Aufgabe{Konvergenzordnung empirisch bestimmen}{6}

Aus den Aufgaben 4 und 5 liegen für jede Methode zwei Fehlerwerte vor: einmal mit Schrittweite $h$ und einmal mit halbierter Schrittweite $h/2$ (Midpoint: $h = 2 \to 1$; Trapez: $h = 2 \to 1$; Simpson: $h = 1 \to 0.5$; die Einzel-Simpson-Regel aus Aufgabe~4(d) --- dort mit $h = b - a = 2$ notiert --- wird hier als zusammengesetzte Regel mit $n = 2$ Teilintervallen und Gitterweite $h = 1$ gelesen).

\begin{enumerate}[label=(\alph*)]
  \item Bilden Sie für jede der drei Methoden den Quotienten $\dfrac{\text{Fehler}(h)}{\text{Fehler}(h/2)}$. \hfill (3~P)
  \item Welche Konvergenzordnung des akkumulierten Fehlers legt jeder Quotient nahe? Vergleichen Sie mit Tabelle~3 des Skripts. (Zur Erinnerung: Bei $O(h^p)$ erwartet man bei Halbierung von $h$ eine Fehlerreduktion um etwa den Faktor $2^p$.) \hfill (2~P)
  \item Sagen Sie ohne Rechnung voraus: Welchen Fehler erwarten Sie ungefähr für die zusammengesetzte Simpson-Regel mit $n = 8$ Intervallen ($h = 0.25$)? \hfill (1~P)
\end{enumerate}

\schreibraum[9]
\Aufgabe{Monte-Carlo-Integration von Hand (Transfer)}{9}

Das Integral $I = \int_0^2 e^x\,dx$ soll nun per Monte Carlo geschätzt werden. Die Integrationsdomäne ist $\Omega = [0, 2]$ mit $|\Omega| = 2$. Gegeben seien $N = 4$ Zufallssamples, gezogen i.i.d.\ aus $\text{Uniform}(0, 2)$:
\[ X_1 = 0.25, \qquad X_2 = 0.75, \qquad X_3 = 1.50, \qquad X_4 = 1.75. \]

\begin{enumerate}[label=(\alph*)]
  \item Berechnen Sie den Monte-Carlo-Schätzer nach Gl.~(46)/(47),
  \[ \hat{I}_N = \frac{|\Omega|}{N}\sum_{i=1}^{N} f(X_i), \]
  und den absoluten Fehler gegenüber dem exakten Wert $\tilde{I}$ aus Aufgabe~4. \hfill (3~P)
  \item Berechnen Sie die erwartungstreue Stichprobenvarianz (unbiased sample variance) nach Gl.~(51),
  \[ \widehat{\text{Var}}(f(\mathbf{X})) = \frac{1}{N-1}\sum_{i=1}^{N}\left(f(X_i) - \bar{f}\right)^2, \]
  mit dem Stichprobenmittel $\bar{f}$, sowie die Stichprobenstandardabweichung $\sigma_f$. Geben Sie alle vier quadrierten Abweichungen an. \hfill (3~P)
  \item Berechnen Sie den Standardfehler (standard error) nach Gl.~(50), $\text{SE}(\hat{I}_N) = \frac{|\Omega|\,\sigma_f}{\sqrt{N}}$. Ist der tatsächliche Fehler aus (a) damit verträglich? \hfill (1~P)
  \item Die Stichprobengröße wird von $N = 4$ auf $4N = 16$ erhöht (bei etwa gleicher Stichprobenvarianz). Um welchen Faktor ändert sich der Standardfehler, und welchen Wert erwarten Sie ungefähr? \hfill (1~P)
  \item Begründen Sie anhand der Formel aus (c), warum der Fehler des Monte-Carlo-Schätzers dimensionsunabhängig (dimension independent) ist --- im Gegensatz zu den Gitterverfahren aus Aufgabe~3. \hfill (1~P)
\end{enumerate}

\schreibraum[12]
\Aufgabe{Lagrange-Interpolation, Simpson-Exaktheit und die SGD-Verbindung (Transfer)}{10}

\begin{enumerate}[label=(\alph*)]
  \item \emph{Von der Interpolation zur Quadratur (Gl.~(45)):} Die Simpson-Regel entsteht, indem $f$ auf $[a, b]$ durch das quadratische Lagrange-Interpolationspolynom
  \[ p(x) = \sum_{j=0}^{2} f(x_j)\,L_j(x), \qquad L_j(x) = \prod_{\substack{k=0 \\ k \neq j}}^{2} \frac{x - x_k}{x_j - x_k}, \]
  durch die drei äquidistanten Knoten $x_0 = a$, $x_1 = m = \frac{a+b}{2}$, $x_2 = b$ ersetzt und dieses exakt integriert wird. Stellen Sie die drei Lagrange-Basispolynome $L_j(x)$ explizit auf (o.\,B.\,d.\,A.\ $a = 0$, $b = h$, $m = \frac{h}{2}$ --- eine Verschiebung des Intervalls ändert die Gewichte nicht) und zeigen Sie durch Integration der Basispolynome, dass sich die Quadraturgewichte
  \[ w_j = \int_0^h L_j(x)\,dx, \qquad w_0 = w_2 = \frac{h}{6}, \quad w_1 = \frac{2h}{3} \]
  ergeben --- also genau die Simpson-Regel $\int_a^b f(x)\,dx \approx \frac{h}{6}\left[f(a) + 4f(m) + f(b)\right]$ aus Aufgabe~4(d). \hfill (3~P)
  \item Zeigen Sie durch explizite Rechnung, dass die Simpson-Regel das Integral $\int_0^1 x^3\,dx$ \emph{exakt} berechnet (ein Intervall, $h = b - a = 1$, Mittelpunkt $m = 0.5$). \hfill (2~P)
  \item Die Simpson-Regel beruht auf einem \emph{quadratischen} Interpolationspolynom durch drei Punkte. Warum ist sie dennoch exakt für alle Polynome bis Grad \emph{drei}? Argumentieren Sie über die Symmetrie des kubischen Restterms um den Intervallmittelpunkt. \hfill (2~P)
  \item \emph{SGD als Monte Carlo:} In Deep Learning ist der wahre Gradient der Mittelwert der Gradienten über alle $n$ Trainingspunkte. Erklären Sie, (i) warum der Mini-Batch-Gradient in SGD ein Monte-Carlo-Schätzer dieses wahren Gradienten ist, (ii) wie die Varianz des Schätzers mit der Batch-Größe $b$ skaliert und welcher Trade-off sich daraus ergibt, und (iii) wie viele SGD-Update-Schritte man im Skript-Beispiel $n = 10^6$, $b = 100$ mit dem Rechenaufwand \emph{eines} Gradient-Descent-Schritts ausführen kann. \hfill (3~P)
\end{enumerate}

\schreibraum[13]
\newpage

\section*{Lösungen --- erst nach Bearbeitung lesen!}

\subsection*{Lösung zu Aufgabe~\ref{auf:1}}

\begin{enumerate}[label=(\alph*)]
  \item Das Integral über ein Intervall ist (nach Division durch die Breite $h$) ein \emph{Mittelwert} der Funktionswerte auf diesem Intervall. Hochfrequente Oszillationen schwanken innerhalb des Intervalls um null herum und mitteln sich heraus (averaging effect), während der niederfrequente Trend erhalten bleibt: Scharfe, schmale lokale Minima (\glqq potholes\grqq) verschwinden, die globale Struktur der Landschaft bleibt sichtbar. Die beiden guten Gründe des Skripts: (1) Integration glättet die Loss-Landschaft und macht Optimierungsmethoden robuster; (2) sie findet tendenziell \emph{robustere Optima}, weil durch den Mittelungseffekt einer Region nicht ein einzelner, zufällig tiefer Punkt, sondern eine im Mittel gute Umgebung bewertet wird.
  \item Die Störkomponente $\sin(2\pi x)$ hat die Periode (Wellenlänge) $T = \frac{2\pi}{2\pi} = 1$. Das Integral einer Sinusschwingung über eine \emph{volle Periode} ist exakt null --- die positiven und negativen Halbwellen interferieren destruktiv. Mit $h = 1$ (oder jedem ganzzahligen Vielfachen) würde der Anteil $\sin(2\pi x)$ also exakt eliminiert:
  \[ \int_{x - 1/2}^{x + 1/2} \sin(2\pi t)\,dt = 0 \quad \text{für jedes } x. \]
  Die Skript-Wahl $h = 0.48 \approx T/2$ liegt nahe an der \emph{halben} Periode: In der Nachbar-Mittelung des Skripts liegen die Auswertungspunkte damit für die schnelle Schwingung nahezu gegenphasig, sodass sich deren Anteile fast vollständig aufheben --- die Kurve in Figure~33 ist daher fast glatt sinusförmig.
  \item Randnotiz S.~65: Die Mittelpunktsregel ist der Spezialfall der Monte-Carlo-Schätzung mit $N = 1$ Sample, das deterministisch genau am Mittelpunkt platziert wird: $\hat{I} = |\Omega| \cdot f(m)$ entspricht $\hat{I}_1 = \frac{|\Omega|}{1} f(X_1)$ mit $X_1 = m$.
\end{enumerate}

\subsection*{Lösung zu Aufgabe~\ref{auf:2}}

\begin{enumerate}[label=(\alph*)]
  \item
  \begin{itemize}[nosep]
    \item \textbf{Mittelpunktsregel:} Rechteck der Höhe $f(m)$ --- konstante Approximation (Polynomgrad 0), ausgewertet am Mittelpunkt.
    \item \textbf{Trapezregel:} Trapez --- lineare Approximation (Grad 1) durch die beiden Endpunkte; Trapeze sind \glqq geometrisch ausdrucksstärker als Rechtecke\grqq{} und erfassen lineare Änderungen.
    \item \textbf{Simpson-Regel:} Parabel --- quadratische Approximation (Grad 2) durch die drei Punkte $a$, $m$, $b$; erfasst Krümmung.
  \end{itemize}
  Die Mittelpunktsregel ist besser als Links-/Rechts-Endpunkt-Regeln, weil der Mittelpunkt die Krümmung der Funktion auf dem Intervall besser approximiert: Der Fehler links vom Mittelpunkt und der Fehler rechts davon haben entgegengesetztes Vorzeichen und heben sich in führender Ordnung auf, während ein Endpunktwert systematisch einseitig daneben liegt.
  \item Zuordnung nach Tabelle~3:
  \begin{center}
  \begin{tabular}{lcc}
  \toprule
  Methode & Einzelintervallfehler & Akkumulierter Fehler \\
  \midrule
  Left/Right & $O(h^2)$ & $O(h)$ \\
  Midpoint & $O(h^3)$ & $O(h^2)$ \\
  Trapezoid & $O(h^3)$ & $O(h^2)$ \\
  Simpson's 1/3 & $O(h^5)$ & $O(h^4)$ \\
  \bottomrule
  \end{tabular}
  \end{center}
  (Der akkumulierte Fehler ist jeweils eine Ordnung niedriger, weil über $n \sim 1/h$ Teilintervalle summiert wird.)
  \item Polynominterpolation hohen Grades oszilliert durch Overfitting --- das ist das Runge-Phänomen: Das Interpolationspolynom schwingt zwischen den Stützstellen auf, mit besonders großen Fehlern an den Intervallrändern. Für Quadraturregeln ab Grad $n \geq 8$ entwickeln die Gewichte zudem \emph{negative} Werte, wodurch die Regeln numerisch instabil werden (große, sich fast auslöschende Beiträge). Daher der praktische Rat \glqq Stop at Simpson\grqq: Höhere Genauigkeit erreicht man robuster über zusammengesetzte Regeln mit mehr Teilintervallen (kleineres $h$) statt über höhere Polynomgrade.
\end{enumerate}

\subsection*{Lösung zu Aufgabe~\ref{auf:3}}

\begin{enumerate}[label=(\alph*)]
  \item Ein Tensorprodukt-Gitter muss jede Kombination der $N$ Stützstellen in jeder der $d$ Dimensionen auswerten --- das sind $N \times N \times \cdots \times N = N^d$ Punkte.
  \begin{itemize}[nosep]
    \item $d = 2$: $100^2 = 10^4$ Punkte --- harmlos.
    \item $d = 10$: $100^{10} = 10^{20}$ Punkte.
  \end{itemize}
  Bei $10^9$ Auswertungen pro Sekunde: $10^{20}/10^9 = 10^{11}$ Sekunden $\approx 3{.}2 \cdot 10^3$ Jahre (rund 3000 Jahre) --- und das für eine \emph{einzige} Integralauswertung, die in iterativer Optimierung tausendfach wiederholt werden müsste.
  \item $N = 2$, $d = 10^6$: $2^{10^6}$ Gitterpunkte. Anzahl der Dezimalstellen: $\log_{10}\!\left(2^{10^6}\right) = 10^6 \cdot \log_{10} 2 \approx 10^6 \cdot 0.30103 \approx 301\,030$, also eine Zahl mit rund $300\,000$ Stellen --- wie im Skript: völlig jenseits jeder Rechenkapazität.
  \item Der Standardfehler von Monte Carlo skaliert mit $O(1/\sqrt{N})$ \emph{unabhängig von der Dimension $d$}: Die Konvergenzrate ist zwar langsamer als bei Quadraturregeln, aber der Aufwand explodiert nicht mit $d$ --- es gibt kein Gitter, sondern nur $N$ zufällige Auswertungspunkte.
\end{enumerate}

\subsection*{Lösung zu Aufgabe~\ref{auf:4}}

\begin{enumerate}[label=(\alph*)]
  \item Exakter Wert:
  \[ \tilde{I} = \int_0^2 e^x\,dx = \left[e^x\right]_0^2 = e^2 - e^0 = 7.389056 - 1 = 6.389056. \]
  \item Mittelpunktsregel, $n = 1$: Mittelpunkt $m = \frac{0 + 2}{2} = 1$, Breite $b - a = 2$:
  \[ \hat{I}_M = 2 \cdot f(1) = 2 \cdot 2.718282 = 5.436564. \]
  Fehler: $|\tilde{I} - \hat{I}_M| = |6.389056 - 5.436564| = 0.952492$.
  \item Trapezregel (Gl.~(44)), $h = b - a = 2$:
  \[ \hat{I}_T = \frac{2}{2}\left[f(0) + f(2)\right] = 1 \cdot (1.000000 + 7.389056) = 8.389056. \]
  Fehler: $|\tilde{I} - \hat{I}_T| = |6.389056 - 8.389056| = 2.000000$.
  \item Simpson-Regel, $h = b - a = 2$, $m = 1$:
  \[ \hat{I}_S = \frac{2}{6}\left[f(0) + 4f(1) + f(2)\right] = \frac{1}{3}\left[1.000000 + 4 \cdot 2.718282 + 7.389056\right] = \frac{1}{3} \cdot 19.262184 = 6.420728. \]
  Fehler: $|\tilde{I} - \hat{I}_S| = |6.389056 - 6.420728| = 0.031672$.
  \item Rangfolge: Simpson ($0.031672$) $<$ Midpoint ($0.952492$) $<$ Trapez ($2.000000$). Bei einer \emph{konvexen} Funktion liegt jede Sekante (Verbindungsgerade der Endpunkte) \emph{oberhalb} der Kurve --- die Trapezfläche überschätzt das Integral. Die Tangente im Mittelpunkt liegt dagegen \emph{unterhalb} der Kurve, und die Mittelpunktsregel entspricht genau der Fläche unter dieser Tangente (der lineare Anteil mittelt sich heraus) --- sie unterschätzt das Integral. Typisch ist zudem, dass der Trapezfehler etwa doppelt so groß ist wie der Mittelpunktfehler, mit umgekehrtem Vorzeichen (hier: $-0.95$ vs.\ $+2.00$).
\end{enumerate}

\subsection*{Lösung zu Aufgabe~\ref{auf:5}}

\begin{enumerate}[label=(\alph*)]
  \item \textbf{Zusammengesetzte Trapezregel} (Gewichtsmuster $1\text{--}2\text{--}2\text{--}\cdots\text{--}2\text{--}1$):
  \[ \int_a^b f(x)\,dx \approx \frac{h}{2}\left[f(x_0) + 2\sum_{i=1}^{n-1} f(x_i) + f(x_n)\right]. \]
  Innere Punkte tragen Gewicht 2, weil jeder innere Punkt zu \emph{zwei} benachbarten Trapezen beiträgt --- einmal als rechter und einmal als linker Endpunkt ---, während die Randpunkte $a$ und $b$ nur zu je einem Trapez gehören.
  \textbf{Zusammengesetzte Simpson-Regel} ($n$ gerade; Gewichtsmuster $1\text{--}4\text{--}2\text{--}4\text{--}\cdots\text{--}4\text{--}1$):
  \[ \int_a^b f(x)\,dx \approx \frac{h}{3}\left[f(x_0) + 4\!\!\sum_{\substack{i=1\\ i\text{ ungerade}}}^{n-1}\!\! f(x_i) + 2\!\!\sum_{\substack{i=2\\ i\text{ gerade}}}^{n-2}\!\! f(x_i) + f(x_n)\right]. \]
  Der Vorfaktor entsteht, weil die einfache Simpson-Regel $\frac{H}{6}[f + 4f + f]$ auf jedes Teilintervall-\emph{Paar} der Breite $H = 2h$ angewandt wird: $\frac{2h}{6} = \frac{h}{3}$ pro Paar [Korrektur ggü.\ Original, S.~64].
  \item Composite Midpoint, $n = 2$, $h = 1$, Mittelpunkte $0.5$ und $1.5$:
  \[ \hat{I}_{M_2} = h\left[f(0.5) + f(1.5)\right] = 1 \cdot (1.648721 + 4.481689) = 6.130410. \]
  Fehler: $|6.389056 - 6.130410| = 0.258646$.
  \item Composite Trapez, $n = 2$, $h = 1$, Gitter $\{0, 1, 2\}$:
  \[ \hat{I}_{T_2} = \frac{1}{2}\left[f(0) + 2f(1) + f(2)\right] = 0.5 \cdot \left[1.000000 + 5.436564 + 7.389056\right] = 0.5 \cdot 13.825620 = 6.912810. \]
  Fehler: $|6.389056 - 6.912810| = 0.523754$.
  \item Composite Simpson, $n = 4$, $h = 0.5$, Gitter $\{0, 0.5, 1, 1.5, 2\}$:
  \begin{align*}
  \hat{I}_{S_4} &= \frac{0.5}{3}\left[f(0) + 4f(0.5) + 2f(1) + 4f(1.5) + f(2)\right] \\
  &= \frac{1}{6}\left[1.000000 + 4 \cdot 1.648721 + 2 \cdot 2.718282 + 4 \cdot 4.481689 + 7.389056\right] \\
  &= \frac{1}{6}\left[1.000000 + 6.594884 + 5.436564 + 17.926756 + 7.389056\right] = \frac{38.347260}{6} = 6.391210.
  \end{align*}
  Fehler: $|6.389056 - 6.391210| = 0.002154$.
\end{enumerate}

\subsection*{Lösung zu Aufgabe~\ref{auf:6}}

\begin{enumerate}[label=(\alph*)]
  \item Fehlerquotienten bei Halbierung der Schrittweite:
  \begin{itemize}[nosep]
    \item Midpoint: $\dfrac{0.952492}{0.258646} \approx 3.68$,
    \item Trapez: $\dfrac{2.000000}{0.523754} \approx 3.82$,
    \item Simpson: $\dfrac{0.031672}{0.002154} \approx 14.7$.
  \end{itemize}
  \item Midpoint und Trapez liegen nahe am Faktor $4 = 2^2$ --- akkumulierter Fehler $O(h^2)$; Simpson liegt nahe am Faktor $16 = 2^4$ --- akkumulierter Fehler $O(h^4)$. Das deckt sich exakt mit Tabelle~3 (Midpoint/Trapezoid: $O(h^2)$, Simpson's 1/3: $O(h^4)$). Dass die empirischen Faktoren leicht unter $4$ bzw.\ $16$ bleiben, liegt daran, dass die Schrittweiten hier noch grob sind und Terme höherer Ordnung mitspielen --- für $h \to 0$ nähern sich die Quotienten den theoretischen Werten.
  \item Für $n = 8$ ($h$ erneut halbiert) erwartet man ungefähr eine weitere Reduktion um den Faktor $16$:
  \[ \text{Fehler}(S_8) \approx \frac{0.002154}{16} \approx 1.3 \cdot 10^{-4}. \]
\end{enumerate}

\subsection*{Lösung zu Aufgabe~\ref{auf:7}}

\begin{enumerate}[label=(\alph*)]
  \item Funktionswerte: $f(X_1) = e^{0.25} = 1.284025$, $f(X_2) = e^{0.75} = 2.117000$, $f(X_3) = e^{1.5} = 4.481689$, $f(X_4) = e^{1.75} = 5.754603$. Summe:
  \[ \sum_{i=1}^4 f(X_i) = 1.284025 + 2.117000 + 4.481689 + 5.754603 = 13.637317. \]
  Schätzer (Gl.~(47)) mit $|\Omega| = 2$, $N = 4$:
  \[ \hat{I}_4 = \frac{2}{4} \cdot 13.637317 = 6.818659. \]
  Fehler: $|\tilde{I} - \hat{I}_4| = |6.389056 - 6.818659| = 0.429603$.
  \item Stichprobenmittel: $\bar{f} = \frac{13.637317}{4} = 3.409329$. Quadrierte Abweichungen:
  \begin{align*}
  (1.284025 - 3.409329)^2 &= (-2.125304)^2 = 4.516917, \\
  (2.117000 - 3.409329)^2 &= (-1.292329)^2 = 1.670114, \\
  (4.481689 - 3.409329)^2 &= (+1.072360)^2 = 1.149956, \\
  (5.754603 - 3.409329)^2 &= (+2.345274)^2 = 5.500310.
  \end{align*}
  Summe: $12.837297$. Erwartungstreue Stichprobenvarianz (Gl.~(51), Nenner $N - 1 = 3$):
  \[ \widehat{\text{Var}}(f(\mathbf{X})) = \frac{12.837297}{3} = 4.279099, \qquad \sigma_f = \sqrt{4.279099} = 2.068598. \]
  \item Standardfehler (Gl.~(50)):
  \[ \text{SE}(\hat{I}_4) = \frac{|\Omega|\,\sigma_f}{\sqrt{N}} = \frac{2 \cdot 2.068598}{\sqrt{4}} = \frac{4.137196}{2} = 2.068598. \]
  Der tatsächliche Fehler $0.429603$ liegt deutlich \emph{innerhalb} einer Standardfehler-Breite --- die Schätzung ist mit ihrer eigenen Unsicherheitsangabe also gut verträglich. ($e^x$ streut auf $[0,2]$ stark, daher ist der SE bei nur $N = 4$ Samples groß.)
  \item Wegen $\text{SE} \propto 1/\sqrt{N}$ ändert sich der Standardfehler bei $N \to 4N$ um den Faktor $\frac{1}{\sqrt{4}} = \frac{1}{2}$ --- er halbiert sich:
  \[ \text{SE}(\hat{I}_{16}) \approx \frac{2.068598}{2} \approx 1.034299. \]
  \item In der Formel $\text{SE} = \frac{|\Omega|\,\sigma_f}{\sqrt{N}}$ kommt die Dimension $d$ \emph{nicht} vor: Der Fehler hängt nur vom Domänenvolumen $|\Omega|$, der Streuung $\sigma_f$ der Funktionswerte und der Samplezahl $N$ ab. Die Konvergenzrate $O(1/\sqrt{N})$ gilt in jeder Dimension gleichermaßen, weil keine Gitterstruktur aufgebaut werden muss --- im Gegensatz zu den $N^d$ Punkten der deterministischen Quadratur aus Aufgabe~3.
\end{enumerate}

\subsection*{Lösung zu Aufgabe~\ref{auf:8}}

\begin{enumerate}[label=(\alph*)]
  \item Mit den Knoten $x_0 = 0$, $x_1 = \frac{h}{2}$, $x_2 = h$ lauten die drei Lagrange-Basispolynome (Gl.~(45)):
  \begin{align*}
  L_0(x) &= \frac{(x - \frac{h}{2})(x - h)}{(0 - \frac{h}{2})(0 - h)} = \frac{2}{h^2}\left(x - \tfrac{h}{2}\right)(x - h), \\
  L_1(x) &= \frac{(x - 0)(x - h)}{(\frac{h}{2} - 0)(\frac{h}{2} - h)} = -\frac{4}{h^2}\,x\,(x - h), \\
  L_2(x) &= \frac{(x - 0)(x - \frac{h}{2})}{(h - 0)(h - \frac{h}{2})} = \frac{2}{h^2}\,x\left(x - \tfrac{h}{2}\right).
  \end{align*}
  (Probe: $L_j(x_k) = 1$ für $j = k$ und $0$ sonst.) Integration der drei Basispolynome über $[0, h]$:
  \begin{align*}
  w_0 &= \frac{2}{h^2}\int_0^h \left(x^2 - \tfrac{3h}{2}x + \tfrac{h^2}{2}\right) dx
       = \frac{2}{h^2}\left[\frac{h^3}{3} - \frac{3h}{2}\cdot\frac{h^2}{2} + \frac{h^2}{2}\cdot h\right]
       = \frac{2}{h^2}\cdot\frac{4h^3 - 9h^3 + 6h^3}{12} = \frac{2}{h^2}\cdot\frac{h^3}{12} = \frac{h}{6}, \\
  w_1 &= -\frac{4}{h^2}\int_0^h \left(x^2 - hx\right) dx
       = -\frac{4}{h^2}\left[\frac{h^3}{3} - \frac{h^3}{2}\right]
       = -\frac{4}{h^2}\cdot\left(-\frac{h^3}{6}\right) = \frac{2h}{3}, \\
  w_2 &= \frac{2}{h^2}\int_0^h \left(x^2 - \tfrac{h}{2}x\right) dx
       = \frac{2}{h^2}\left[\frac{h^3}{3} - \frac{h}{2}\cdot\frac{h^2}{2}\right]
       = \frac{2}{h^2}\cdot\frac{h^3}{12} = \frac{h}{6}.
  \end{align*}
  Damit gilt
  \[ \int_a^b f(x)\,dx \approx \int_a^b p(x)\,dx = \frac{h}{6}\,f(a) + \frac{2h}{3}\,f(m) + \frac{h}{6}\,f(b) = \frac{h}{6}\left[f(a) + 4f(m) + f(b)\right], \]
  also genau die Simpson-Regel aus Aufgabe~4(d). Plausibilitätsprobe: $w_0 + w_1 + w_2 = \frac{h}{6} + \frac{2h}{3} + \frac{h}{6} = h$ --- die konstante Funktion $f \equiv 1$ wird exakt integriert.
  \item Exakter Wert: $\int_0^1 x^3\,dx = \left[\frac{x^4}{4}\right]_0^1 = \frac{1}{4} = 0.25$. Simpson mit $h = 1$, $m = 0.5$:
  \[ \hat{I}_S = \frac{1}{6}\left[f(0) + 4f(0.5) + f(1)\right] = \frac{1}{6}\left[0 + 4 \cdot 0.125 + 1\right] = \frac{1}{6} \cdot 1.5 = 0.25. \]
  Die Simpson-Regel liefert also den exakten Wert $\frac{1}{4}$ --- Fehler null.
  \item Per Konstruktion --- Integration des quadratischen Lagrange-Interpolanten durch drei Punkte, vgl.\ Teilaufgabe~(a) --- ist Simpson exakt für Polynome bis Grad 2. Für ein kubisches Polynom mit Leitkoeffizient $c$ betrachte man die Abweichung vom quadratischen Interpolanten: Da sie ein kubisches Polynom mit Nullstellen an den drei Knoten $a, m, b$ ist, gilt exakt $c\,(x-a)(x-m)(x-b) = c\left[(x-m)^3 - \tfrac{h^2}{4}(x-m)\right]$ --- beide Summanden und damit die gesamte Abweichung sind \emph{ungerade} Funktionen bezüglich des Intervallmittelpunkts $m$. Ihr Integral über das (um $m$ symmetrische) Intervall verschwindet daher exakt --- der negative Beitrag links von $m$ hebt den positiven rechts von $m$ auf. Deshalb integriert die Simpson-Regel auch jeden kubischen Anteil exakt und ist insgesamt exakt für alle Polynome bis Grad 3 --- ein \glqq Genauigkeitsgrad geschenkt\grqq. Für $\sin(x)$ gilt das nicht: Dort bleibt ein kleiner Restfehler [vgl.\ Korrektur ggü.\ Original, S.~67: $S_2 \approx -0.956788$ mit Fehler $\approx 3.4 \cdot 10^{-4}$, nicht exakt null].
  \item
  \begin{itemize}[nosep]
    \item[(i)] Der wahre Gradient ist ein Mittelwert über alle $n$ Trainingspunkte --- formal ein Integral bzw.\ eine Erwartung über die Datenverteilung. Der Mini-Batch-Gradient mittelt die Gradienten über eine \emph{zufällig gezogene} Teilmenge (Batch) der Größe $b$: Er ist genau der Monte-Carlo-Schätzer aus Gl.~(46)/(47) für diese Erwartung --- erwartungstreu (Gl.~(48)), aber verrauscht.
    \item[(ii)] Analog zu Gl.~(49) skaliert die Varianz des Schätzers ungefähr mit $1/b$ (Standardfehler $\propto 1/\sqrt{b}$). Trade-off: Große Batches geben rauscharme, genaue Gradienten, kosten aber mehr Rechenzeit pro Update und nehmen das Gradientenrauschen weg, das beim Entkommen aus scharfen lokalen Minima hilft; kleine Batches sind billig und verrauscht --- das Rauschen wirkt wie eine eingebaute Glättung/Exploration der Loss-Landschaft. Die Batch-Größe steuert also direkt den Noise--Varianz-Trade-off des Trainings.
    \item[(iii)] Ein voller Gradient-Descent-Schritt kostet $n = 10^6$ Gradientenauswertungen, ein SGD-Schritt nur $b = 100$. Pro Update ist SGD also $n/b = 10^4$-mal billiger: Mit dem Rechenaufwand \emph{eines} GD-Schritts führt SGD $10\,000$ Update-Schritte aus.
  \end{itemize}
\end{enumerate}

\bigskip
\hrule
\medskip
\noindent\small\textbf{Hinweis zu Korrekturen:} Dieses Übungsblatt verwendet durchgehend die gegenüber dem Original-PDF korrigierten Werte (vgl.\ Errata-Liste): Composite-Simpson-Vorfaktor $h/3$ pro Teilintervall-Paar [S.~64]; Midpoint $n{=}2$ am Skript-Beispiel $\int_{-2}^{-1}\sin(x)\,dx$: $\hat{I} \approx -0.966485$, Fehler $\approx 0.010036$ [S.~66]; Trapez $T_2 \approx -0.93644$, Fehler $\approx 0.02001$ [S.~67]; Simpson $S_2 \approx -0.956788$, Fehler $\approx 3.4 \cdot 10^{-4}$ --- Simpson ist exakt für Polynome bis Grad 3, nicht für $\sin(x)$ [S.~67].

\end{document}
